\documentclass[pstricks,border=12pt]{standalone}
\usepackage{amsmath,amsthm,array,ifthen,amsfonts}
\usepackage{amsmath,pst-coil,bigints}
\usepackage{pstricks-add,pst-plot,pst-node,pst-3dplot}
\usepackage{bm,calculus}
\psset{showorigin=false}
\begin{document}
\begin{pspicture}(-2.5,-0.8)(2,3.1) % \psgrid
\psset{algebraic,Alpha=80,Beta=20,plotstyle=curve,xunit=0.7,yunit=0.8}
\pstThreeDCoor[linecolor=black,Dy=2,deltay=2,xMin=-4.3,xMax=3,yMin=-3,yMax=3.1,zMin=-1,zMax=4.2,spotZ=180,IIIDticks,IIIDlabels]
\pscustom[fillstyle=solid,fillcolor=gray!60,opacity=0.75,linestyle=none]{
\parametricplotThreeD(0,\psPiTwo){(2-2*cos(t))*cos(t)|(2-2*cos(t))*sin(t)|0}
} % sombra en la base
\parametricplotThreeD[linecolor=black!60,linewidth=0.8pt](0,\psPiTwo){2*(1-cos(t))*cos(t)|2*(1-cos(t))*sin(t),0} % curva base
\pscustom[fillstyle=vlines,hatchcolor=gray!60,hatchwidth=0.6pt,hatchsep=1pt,linestyle=none]{
\parametricplotThreeD(0,\psPiTwo){(2-2*cos(t))*cos(t)|(2-2*cos(t))*sin(t)|1-((2-2*cos(t))*sin(t))/3}
} % sombra en la base

%\parametricplotThreeD[linecolor=black!60,linewidth=0.8pt](0,\psPiTwo){2*(1-cos(t))*cos(t)|2*(1-cos(t))*sin(t),1-(2*(1-cos(t))*sin(t))/3} % curva base

\parametricplotThreeD[linecolor=black,linewidth=0.2pt](0,\psPiTwo)(0,1){2*(1-cos(t))*cos(t)|2*(1-cos(t))*sin(t),u*(3-2*sin(t)+2*cos(t)*sin(t))/3} % cilindro
\parametricplotThreeD[linecolor=black,linewidth=0.2pt](0,1)(0,\psPiTwo){2*(1-cos(u))*cos(u)|2*(1-cos(u))*sin(u),t*(3-2*sin(u)+2*cos(u)*sin(u))/3} % cilindro
\pstThreeDPut(-2,1.5,0.2){$Q$}
\pstThreeDPut(-2,1,0){$\bullet$}
\pstThreeDLine[linecolor=black!40,linewidth=1.2pt](-2,1,0)(-2,1,0.5)
\pscustom[fillstyle=hlines,hatchcolor=gray!30,hatchwidth=0.8pt,hatchsep=0.8pt,linestyle=none]{
\parametricplotThreeD(0,\psPiTwo){(2-2*cos(t))*cos(t)|(2-2*cos(t))*sin(t)|1-((2-2*cos(t))*sin(t))/3}
} 
\parametricplotThreeD[linecolor=black!40,linewidth=1pt](0,\psPiTwo){2*(1-cos(t))*cos(t)|2*(1-cos(t))*sin(t),1-(2*(1-cos(t))*sin(t))/3} 
\pstThreeDLine[linecolor=black!50,linewidth=0.3pt](0,0,0)(0,0,4)
\pstThreeDPut(-2,1,0.5){$\bullet$}
\pstThreeDPut(-2,1.5,0.7){$P$}
\end{pspicture}

\end{document}
